%!TEX root = ./main.tex

\section[5G Core]{The Core Becomes Software}

% SECTION TIMING: 8:00--30:00 (4 frames). This is the first of Wednesday's two
% promised repeats of Monday's rule: the core splits again. The CU/DU/RU frame at
% the end is the second. Both are set up by "You have now watched the same move
% twice" at the end of Monday -- point back to it explicitly.


\begin{frame}{5G made the core look more like a cloud platform}
  \begin{columns}[T,onlytextwidth]
    \column{0.56\textwidth}
      \centering
      \begin{tikzpicture}
        \node[corenode,minimum width=1.25cm,font=\scriptsize\bfseries] (amf) at (0,1.25) {AMF};
        \node[corenode,minimum width=1.25cm,font=\scriptsize\bfseries] (smf) at (1.8,1.25) {SMF};
        \node[corenode,minimum width=1.25cm,font=\scriptsize\bfseries] (ausf) at (3.6,1.25) {AUSF};
        \node[corenode,minimum width=1.25cm,font=\scriptsize\bfseries] (udm) at (0,0.15) {UDM};
        \node[corenode,minimum width=1.25cm,font=\scriptsize\bfseries] (pcf) at (1.8,0.15) {PCF};
        \node[corenode,minimum width=1.25cm,font=\scriptsize\bfseries] (nrf) at (3.6,0.15) {NRF};
        \node[oranode,minimum width=5.35cm,minimum height=0.58cm,font=\scriptsize\bfseries] (bus) at (1.8,-0.95) {Service-based interfaces};
        \foreach \n in {amf,smf,ausf,udm,pcf,nrf} \draw[thin,draw=coreorange] (\n) -- (bus);
        \node[netnode,minimum width=1.35cm,font=\scriptsize\bfseries] (ran) at (-0.05,-2.25) {RAN};
        \node[corenode,minimum width=1.35cm,font=\scriptsize\bfseries] (upf) at (1.8,-2.25) {UPF};
        \node[netnode,minimum width=1.55cm,font=\scriptsize\bfseries] (data) at (3.65,-2.25) {Data net};
        \draw[flow] (ran)--(upf); \draw[flow] (upf)--(data);
      \end{tikzpicture}

      \vspace{0.3em}
      {\scriptsize The SMF steers the UPF over the N4 interface.\par}
      \source{https://www.nokia.com/asset/i/210365/}{Nokia 5G Core Service Based Architecture}
    \column{0.40\textwidth}
      \small
      \plainterm{AMF}{Access + Mobility}\par\vspace{0.35em}
      \plainterm{SMF}{Session Management}\par\vspace{0.35em}
      \plainterm{UPF}{User Plane Function}\par\vspace{0.35em}
      \plainterm{AUSF/UDM}{identity}\par\vspace{0.35em}
      \plainterm{PCF}{policy}\par\vspace{0.35em}
      \plainterm{NRF}{service registry}
  \end{columns}

  \vspace{0.3em}
  \takeaway{Core functions now discover and call each other as services.}
\end{frame}

% Teaching note (90 s): this is the third \gonebox and students should be ahead of
% you by now. Ask them to predict the seam before you show it. That prediction is
% the whole point of having built the pattern on Monday.
\begin{frame}{Third time. Where did the MME go?}
  \gonebox
    {In 4G the MME was the one box that made control decisions: who may
     attach, where you are, which gateway serves you.

     \vspace{0.4em}
     One box, every kind of decision.}
    {Those decisions do not resemble each other.

     \vspace{0.4em}
     Checking who you are happens once. Tracking where you are happens
     constantly. Setting up a session happens per app.

     \vspace{0.4em}
     Same box, three completely different loads.}
    {5G cut the control plane itself.

     \vspace{0.4em}
     AMF handles access and mobility. SMF handles sessions. Each is a
     separate service that scales on its own.}

  \vspace{0.5em}
  \glossline{AMF = Access and Mobility Management Function \quad
    SMF = Session Management Function}

  \takeaway{Same rule, one level deeper: now the control plane splits.}
\end{frame}

\begin{frame}{A 5G radio does not guarantee a 5G core}
  \vspace{-0.2em}
  {\small NSA = Non-Standalone \quad SA = Standalone \quad
    QoS = Quality of Service\par}
  \vspace{0.3em}
  \begin{columns}[T,onlytextwidth]
    \column{0.48\textwidth}
      \centering
      \Large\textbf{NSA}

      \vspace{0.6em}
      \begin{tikzpicture}[node distance=0.35cm,scale=0.86,transform shape]
        \node[netnode] (lte) {LTE anchor};
        \node[netnode,below=of lte] (nr) {5G NR};
        \node[corenode,right=of lte] (epc) {4G EPC};
        \draw[flow] (lte) -- (epc);
        \draw[flow] (nr) -- (lte);
      \end{tikzpicture}

      \vspace{0.6em}
      \small The phone shows 5G.\\Control still runs through\\LTE and the 4G core.
    \column{0.48\textwidth}
      \centering
      \Large\textbf{SA}

      \vspace{0.6em}
      \begin{tikzpicture}[node distance=0.35cm,scale=0.86,transform shape]
        \node[netnode] (nr) {5G NR};
        \node[corenode,right=of nr] (gc) {5G Core};
        \draw[flow] (nr) -- (gc);
      \end{tikzpicture}

      \vspace{1.35em}
      \small A 5G radio on a 5G core.\\Slicing, edge, and new QoS\\become possible.
  \end{columns}

  \vspace{0.6em}
  \takeaway{The icon reports the radio.\\The features come from the core.}
  \source{https://about.att.com/blogs/2025/5g-standalone-nationwide.html}{AT\&T: ``Not all 5G is created equal,'' 2025}
\end{frame}

\begin{frame}{The 5G base station can be split across locations}
  \vspace{0.3em}
  \centering
  \begin{tikzpicture}[node distance=1.15cm,scale=0.80,transform shape]
    \node[netnode,minimum width=1.30cm] (ue) {UE};
    \node[netnode,minimum width=1.70cm,right=of ue,font=\scriptsize\bfseries] (ru) {RU};
    \node[netnode,minimum width=1.70cm,right=of ru,font=\scriptsize\bfseries] (du) {DU};
    \node[netnode,minimum width=1.70cm,right=of du,font=\scriptsize\bfseries] (cu) {CU};
    \node[corenode,minimum width=1.70cm,right=of cu,font=\scriptsize\bfseries] (core) {5G Core};
    \draw[flow] (ue)--node[above,font=\tiny]{radio}(ru);
    \draw[flow] (ru)--node[above,font=\tiny]{fronthaul}(du);
    \draw[flow] (du)--node[above,font=\tiny]{F1}(cu);
    \draw[flow] (cu)--node[above,font=\tiny]{N2/N3}(core);
    \node[font=\scriptsize,below=0.35cm of ru,align=center] {RF +\\lower PHY};
    \node[font=\scriptsize,below=0.35cm of du,align=center] {upper PHY\\MAC / RLC};
    \node[font=\scriptsize,below=0.35cm of cu,align=center] {PDCP\\RRC / SDAP};
  \end{tikzpicture}

  \vspace{0.7em}
  {\small One gNodeB can be a system spread across a tower, a cabinet, and a data center.\par}

  \glossline{RU / DU / CU = Radio / Distributed / Central Unit\\
    PHY = physical layer \quad MAC = Medium Access Control \quad
    RLC = Radio Link Control\\
    PDCP = Packet Data Convergence Protocol \quad RRC = Radio Resource Control\\
    SDAP = Service Data Adaptation Protocol}

  \qa{Who owns each box, and who owns the interfaces between them?}
     {In a classic RAN, one vendor owns all of it. Splitting the base station only
      helps if the seams between the pieces are specified well enough for a second
      vendor to plug in.}
\end{frame}
